\scanpage{76}
% lines 1–2, heading and superscribed qualification
$P=P_1*P_2$ \inserted{($P_i$ torseur $\unsure{\chi|\gamma}$-trivialisé, sous module $L_i$)} comme solution d'un pb universel
% lines 3–8, objects of the problem
On s'intéresse aux torseurs $\Pi$ sous un module $M$, et aux applications
\[
P_1\times P_2\xrightarrow{f}\Pi
\]
(qu'on va appeler ``biaffines'' compatibles avec les $\chi$-trivialisations données $\pi_1,\pi_2$ de $P_1,P_2$) qui satisfont la condition suivante :
% lines 9–14, bilinear datum and translations
\[
\exists\quad L_1\times L_2\xrightarrow{\alpha}\deleted{\unsure{\Pi}}M,\quad\text{application bilinéaire,}\tag*{$(*)$}
\]
telle qu'on ait
\[
\left\{\begin{aligned}
\deleted{\alpha}f(x_1+u_1,x_2)&=\alpha(x_1,x_2)+\deleted{B}\alpha(u_1,\pi_2(x_2))\\
\deleted{\alpha}f(x_1,x_2+u_2)&=\alpha(x_1,x_2)+\deleted{B}\alpha(\pi_1(x_1),u_2)
\end{aligned}\right.
\]
(d'où
\[
\begin{aligned}
\deleted{\alpha}f(x_1+u_1,x_2+u_2)={}&\alpha(x_1,x_2)+\deleted{B}\alpha(u_1,\pi_2(x_2))+
\deleted{B}\alpha(\pi_1(x_1),\unsure{x_2|u_2})\\
&+\chi\,\deleted{B}\alpha(u_1,u_2)
\end{aligned}
\]
% lines 15–21, dependence and qualification about uniqueness
Cela signifie donc que $\alpha$ est connu \inserted{(\unsure{en termes de|au lieu de|uniquement de} $B$)} quand on connait sa valeur sur \uline{un} couple $(x_1,x_2)$ de sections, et qu'il est donné par la formule précédente.

\deleted{On trouve que $B$ est} En fait, \deleted{$B$}$\alpha$ n'est pas néc. unique :
(car $(E_1\otimes L_2)\oplus(L_1\otimes E_2)\longrightarrow L_1\otimes L_2$ n'est pas néc. surjectif), et on va définir un \deleted{\unsure{Bi}} ``morphisme biaffine'' comme un couple $(f,\alpha)$, satisfaisant la condition de variance ci-dessus.
% archive sheet number: 75
